\section{What the literature contributes}
\label{sec:literature}

Computational law has developed several answers to different parts of this
problem. Logic programming makes statutory dependencies executable. Languages
such as Catala organize definitions and exceptions close to legal text. Normative
languages represent duties, permissions and violations. Contract languages model
interactions through time. Recent language-model research attempts to cross the
remaining boundary from prose to these representations. A useful product combines
these contributions selectively, because their guarantees concern different
objects.

\subsection{Scope and method of this review}

The review began with the seven requested links and expanded through their
references, forward-citation searches for Catala and Stipula, and searches on
executable law, legal reasoning, formalization, dates, testing and process
compliance. ResearchAssistant supplied discovery records and local PDF extraction.
The saved OpenAlex citation searches returned 52 entries for the selected Catala
record and 13 for the selected Stipula journal record. Those are retrieval counts
for particular indexed editions, not complete or comparable measures of influence.
Preprints, revised editions and new publications can be missing or duplicated.

The local library contains 25 paper editions representing 24 distinct works.
The AAAI tax-reasoning paper and its extended preprint are retained separately;
all seven requested sources are present. Technical reading covered the relevant
methods, semantics, evaluation and supporting appendix material. The accompanying
notes identify the inspected sections, version differences, implementation checks
and limits of adoption \citep{libraryreview}. This is a focused technical review
supporting a product decision, rather than a claim to have enumerated every
citation. Unavailable secondary leads are recorded separately and do not support
implementation claims.

Version control is particularly important here. The ARc paper was first posted
in 2025, but the inspected second version is from July 2026. The eFLINT paper's
inspected version is from August 2026. The Stipula preprint requested by the
sponsor has three authors; the later journal edition has a different author list.
Bibliographic records and local filenames make these distinctions explicit.
Published results are reported as their authors' findings; none of their
benchmarks or proof developments was independently rerun for this proposal.

\subsection{From statutory text to a language of definitions and exceptions}

The British Nationality Act logic program is an early demonstration that selected
legal provisions can be represented as executable dependencies. Its significance
for this project is architectural: a conclusion can be explained through the
conditions supporting it, while rules and facts are represented separately.
The work also confronts negative information and temporal conditions, which
remain central to bank data. It does not show that arbitrary legal prose has a
unique automatic translation \citep{sergot1986}.

Catala addresses a more specific obstacle: legal definitions often state a
general rule and then exceptions to that rule. Its default calculus represents
exceptions explicitly, distinguishes an absent applicable value from a conflict,
and gives those situations operational meaning. In the core construction,
multiple applicable exceptions can produce a conflict even when their computed
values happen to agree. That is a deliberate semantics, not an invitation to
resolve legal priorities by file order \citep[secs.~3--4]{catala2021}.

For LegalMath, this suggests a representation in which an exception names the
rule it qualifies and preserves the source supporting its priority. A discount
exception to a restriction on gifts should remain visibly attached to that
restriction. Two independently imported exceptions that both apply should trigger
a review of their relationship, not an arbitrary choice by the language model.
The product should also distinguish this normative priority from a document's
publication date: a newer explanatory document does not mechanically override
every older source it mentions.

The Catala compilation result is relevant but narrower than the phrase
``verified legal software'' might suggest. The paper proves properties of a
translation from the core default calculus into a lambda calculus with exceptions,
with a mechanized development in F*. It distinguishes that model from the
separately implemented production compiler. The inspected repository's
formalization notes additionally disclose an admitted supporting lemma; this
review has not replayed the development \citep[sec.~4]{catala2021,catalacode}.
The useful precedent is to define a small semantics and state exactly which
translation preserves it. It provides no theorem that a human or model has
captured the meaning of an SFC circular.

Catala's literate presentation also addresses the organizational problem.
\citet[sec.~4.2]{huttner2022} describe collaboration between legal experts and
programmers around a common representation. The legal expert can challenge a
condition without reverse-engineering application code, and the engineer can
make an ambiguity concrete. The exploratory lawyer study in the original Catala
paper gives useful interface evidence, but its small educational setting does
not establish reliable defect detection by bank reviewers
\citep[sec.~6]{catala2021}. LegalMath should borrow the collaboration model and
test its own interface with the intended users.

Controlled language offers another way to make formal statements reviewable.
Logical English retains a readable rule syntax while assigning a computational
meaning to predicates and variables \citep{logicalenglish2023}. For this product,
a controlled paraphrase should be generated from the adopted rule, so changing
the rule changes its explanation. It should display how variables are bound:
``the client'' must refer to the same person throughout a condition. Readability
cannot substitute for a defined semantics, and unconstrained English should not
be presented as if it were executable merely because it looks regular.

\subsection{Dates and test generation expose defects that ordinary examples miss}

The date-arithmetic work of \citet{dates2024} demonstrates why even a simple phrase
such as adding a month needs a convention. Calendar months have different
lengths; leap days create boundary cases; sequences of date operations need not
behave like addition of integers. The paper develops explicit rounding choices,
formalizes date semantics and selected properties, and uses static analysis to
detect ambiguities in legal computations. Its application to benefit rules
connects the formalism to realistic source language rather than treating dates
as a generic programming-library problem.

The immediate lesson for annual reviews and three-year transaction windows is
to preserve the chosen calendar operation and its source or interpretation.
``Three years'' must not silently become a fixed number of days. A leap-day case
can expose disagreement between two plausible implementations. The paper's
day-level formalization does not settle intraday questions such as withdrawal
instructions arriving at a trading cut-off, nor does the mechanization make the
whole static-analysis implementation verified
\citep[secs.~2--6]{dates2024}. Those are distinct prototype obligations.

CUTECat complements this work by generating test cases from Catala programs using
concolic execution: an execution carries concrete values and symbolic conditions,
then a solver seeks inputs that take a different branch. The method specializes
its exploration to default expressions and uses preferences to obtain inputs
that are easier for humans to understand. In the reported evaluation, generated
cases exposed 20 manually validated seeded mutations. A much larger example also
shows that exhaustive-looking exploration can produce a substantial volume of
cases and runtime \citep[secs.~3--5]{cutecat2025}.

LegalMath should borrow branch-directed and mutation-based testing, especially
for inclusive boundaries, alternatives, exceptions and rounding. The result is
evidence about the encoded program. A branch corresponding to an omitted
paragraph does not exist for the test generator to discover. Independent
source-derived scenarios are therefore essential. The official CUTECat artifact
documentation was inspected, but its large execution image was not run; the
prototype should first implement a small set of useful mutations before adopting
that full toolchain \citep{libraryreview}.

\subsection{Legal rules need authority, unknown facts and duties}

LegalRuleML is relevant because it treats legal meaning as more than a Boolean
expression. The standard addresses modalities, defeasibility, temporal aspects,
sources and contextual information. Earlier work on legal interpretations shows
how alternative readings and their contexts can be represented rather than
collapsed prematurely into one rule
\citep[sec.~4]{legalinterpretations2014,legalruleml2021}. This is a strong basis
for LegalMath's metadata: every executable condition should identify its source,
scope, adopted interpretation and relationship to other rules.

The standard is an interchange representation, however. XML that names an
obligation does not by itself define the event processor that creates, satisfies
or breaches it. LegalMath should adopt a deliberately small execution profile
with explicit meaning, and use LegalRuleML concepts for interoperability. Requiring
compliance officers to review raw XML would not solve the bank's usability problem.

s(LAW) illustrates a different benefit of formal legal reasoning: explaining
conclusions under alternative factual assumptions. Its use of s(CASP) combines
logic-programming rules, explicit negative information and default negation, and
produces justifications. The school-admission examples demonstrate discretionary
conditions and hypothetical cases, including assumptions that affect which
conclusions can be derived \citep[secs.~4--5]{slaw2024}.

This distinction matters when an SPI record lacks an agreement or a qualification.
A hypothetical answer can say which additional fact would establish a condition.
It cannot treat that fact as already established. LegalMath should display
``would hold if this evidence were confirmed'' separately from an evaluated
client fact. Different treatments of negation must also remain explicit: absence
from a database, a documented negative statement and a defeasible assumption are
three different things. The s(LAW) examples support explanation and exploration;
the linked implementation could not be fully audited in this review.

eFLINT contributes an especially useful operational distinction. A prohibited or
otherwise impermissible action can still occur in the real world and have legal
effects; the system must be able to record the resulting violation. A pure
authorization engine that simply removes such transitions cannot audit what
actually happened. The language's duties, actions and open types are relevant
to missing evidence and ongoing controls \citep[secs.~3--5]{eflint2026}.

The 2026 eFLINT account also describes differences between implementations and
versions. It notes the absence of a formal operational semantics and type system
for version 4, and limitations concerning fine-grained explanations and normative
priority; another implementation has a stable-model foundation. Guarantees must
therefore be attributed to the actual version and engine used
\citep[secs.~6--7]{eflint2026}. The prototype should borrow the distinction between
permission, occurrence and violation without assuming that installing an eFLINT
interpreter supplies all the required proof evidence.

\subsection{Stipula and the transition from a decision to a contract history}

Stipula makes parties, agreement, state, assets, permitted functions and scheduled
events explicit. Its semantics describes how a contract evolves when parties act
and time advances. This is well aligned with consent, withdrawals, review events
and requests for evidence. It also provides a notion of legal bisimulation,
relating contracts through their observable interactions
\citep[secs.~3--5]{stipula2021}.

The important borrowing is the state model, not a presumption that bank compliance
requires a blockchain. The original paper discusses centralized as well as
distributed execution. A bank can record a source-linked event history in its
ordinary systems. Moreover, a similarity of formal interaction does not establish
equivalence of the source texts as a matter of law. The accompanying discussion
of legal calculi highlights how executable arrangements interact with external
intervention and legal choices that code can otherwise suppress
\citep{stipula2021,crafa2022}.

Several semantic choices need explicit review. In the studied Stipula semantics,
scheduled events take precedence over function execution at the same time. That
is a clear rule for the language; it is not an automatic answer to the bank's
cut-off or consent-ordering questions. The prototype should construct a trace in
which a withdrawal and a transaction have the same recorded time, and identify
the finer ordering evidence or adopted policy required to decide it.

The Stipula-to-Java/JML work offers a concrete verification route. It translates
contract structure into Java together with specifications and applies KeY to
properties of the result. Four contract examples establish feasibility for the
reported translation and properties. The inspected implementation includes a
translator and generated examples \citep[secs.~3--4 and apps.~A--B]{stipulakey2025,stipulakeycode}.
This is worth a bounded experiment if workflow verification proves important.

The encoding also defines the limit of the evidence. Event conditions are
abstracted through symbolic Boolean values rather than a complete implementation
of banking calendars. Automated handling of the relevant loops uses inputs that
remain constant across iterations. A proof about that model does not establish
correctness for arbitrary changing portfolios, unbounded event queues or all
real-time orderings. LegalMath must state and test any adapter's correspondence
to its own event semantics before using the proof in a compliance package.

The 2026 reachability paper supplies a direct design constraint. Through
counter-machine reductions it establishes undecidability for general and several
restricted microStipula fragments. It also identifies a decidable fragment,
called DI, with both disjoint function/event source states and immediate events,
using a well-structured transition-system argument
\citep[secs.~2--4]{reachability2026}. DI is not the same condition as the
disjoint-cycle construction in the Java/JML paper.

The product implication is to ask bounded, named questions. For example, can an
approved rule set permit streamlined treatment after a recorded withdrawal in
any trace of at most $k$ events under a specified ordering model? A counterexample
is useful immediately. Failure to find one supports only the checked bound unless
an additional invariant proof establishes a broader claim. A timeout or unknown
result must remain unresolved. The research does not imply that a small acyclic
financial decision is intractable; it limits the promise of a universal verifier
for unrestricted event-based contracts.

\subsection{Consistency of clauses and compliance of processes}

ContractCheck models legal contracts through structured clauses, an ontology and
constraints, and uses automated reasoning to identify consistency and execution
problems. Its share-purchase-agreement case study is particularly relevant to
the presentation of contradictions: reviewers need to know which claims cannot
coexist and which circumstances make an obligation executable
\citep[secs.~3--7]{contractcheck2025}.

Two distinctions are essential when borrowing this approach. First, the study's
input formalization is manual; its reasoning results do not measure automatic
translation fidelity. Second, showing that each claim can be executed under
some circumstances differs from showing that all duties can be satisfied in one
execution, and both differ from showing that every possible execution complies.
LegalMath should label each of those questions separately. If a solver returns
an unsatisfiable core, that identifies conflicting constraints; a smallest or
minimal explanation requires a further minimization procedure rather than an
assumption about the core \citep[secs.~4--5, 8]{contractcheck2025}.

Process-compliance research reinforces the need to connect rules with activities,
actors and the effects of execution. \citet{ghose2007} describe effect accumulation,
constraint checking and repair in business processes. For the bank, the analogue
is the difference between determining that a warning is required and showing
that the relevant warning was delivered before the required event.

\citet[sec.~2]{compliancehard2015} distinguish achievement obligations, which must
be accomplished, maintenance obligations, which must hold across an interval,
and punctual obligations. They also distinguish partial compliance, where some
execution complies, from full compliance, where all executions do. Their
complexity results concern specified process and obligation languages
\citep[secs.~3--4]{compliancehard2015}. The prototype should borrow those questions
and delimit its model rather than import a complexity label for all regulation.
The practical consequence is that a flowchart showing one successful path is
insufficient evidence about the paths that omit a review or mishandle an exception.

\subsection[Language models and the formalization boundary]{Language models can propose formalizations; their inputs remain the weak point}

The SARA tax-reasoning dataset makes explicit the ingredients an evaluation needs:
statutory text, cases, formal rules and expected conclusions. Its corpus uses
selected and simplified US tax provisions and constructed examples. It is a
valuable experimental design, but neither a current tax authority nor a test set
representative of SFC circulars \citep[secs.~2--3]{sara2020}.

The later work connecting statutory reasoning to information extraction associates
facts with spans in case descriptions and evaluates their effect on reasoning.
This is a useful model for recording where a client attribute came from, which
person it describes and how it was normalized
\citep[secs.~3--5]{connecting2023}. Its task-specific date imputations and
semi-synthetic construction require caution: an extraction convention suitable
for that dataset is not an evidential rule for missing bank records
\citep[sec.~7 and apps.~B, E]{connecting2023}.

ContractNLI adds a complementary annotation pattern. It separates supported,
contradicted and not-mentioned hypotheses, with evidence spans in contracts.
Its 607 non-disclosure agreements and fixed hypotheses expose the importance of
exceptions and evidence spread across a document. It supplies useful ideas for
an evidence-review interface, while evaluating document inference rather than
compilation of regulatory rules \citep[secs.~2--5]{contractnli2021}.

Logic-LM and LINC both divide language interpretation from symbolic inference.
Logic-LM routes formalized problems to solvers and uses error feedback to refine
malformed programs. LINC translates into first-order logic and uses a theorem
prover, with multiple samples and voting in its experiments
\citep[sec.~3]{logiclm2023}\citep[secs.~2--5]{linc2023}. The transferable idea is a
clear boundary: the model proposes the statement, and the solver reasons over
that statement. Their benchmarks do not establish legal translation fidelity.
Successful parsing can leave a wrong quantifier untouched, and repeated models
can agree on the same missing exception.

The inspected Logic-LM repository also offers fallback strategies when formal
execution fails, including a language-model answer or random selection for its
benchmark setting. Those are not acceptable defaults for this prototype's
regulatory decisions. A failed formalization should produce an unresolved issue
with a reason. Research code can be valuable while requiring a different
operational policy \citep{libraryreview}.

The ARc paper is especially close to the proposed product. It separates creation
and human vetting of reusable policy models from subsequent verification of
natural-language material. Its workflow uses source passages, variable
identification, repeated translations and deterministic explanations. The
verification algorithm checks consistency and relationships between the policy,
premises and claim, and distinguishes invalidity, validity, satisfiability and
several forms of inability to decide
\citep[sec.~3 and Algorithm~1]{arc2025}.

This is a strong precedent for an interpretation workspace with explicit
abstention. It is also a warning about terminology. The paper's reported
``soundness'' metric is $1-\mathrm{FP}/N$, an empirical rate over its decisions;
it is not a theorem establishing the translator's semantic soundness. At the
3/3 threshold in Table 1, 10 false positives among 1,689 decisions yield about
99.4 percent on that metric, while reported precision is 94.4 percent and valid
recall is 14.9 percent. The 563 items were each run three times, so the total is
not 1,689 independent policies \citep[sec.~4, Table~1]{arc2025}.

Those results motivate measuring both errors and abstention. A cautious system
may produce few confident errors by deciding very little. Its business value
then depends on whether the unresolved cases become easier to review. The ARc
paper also discusses difficulties with cross-references, tables and implicit
assumptions, all present in the bank's problem
\citep[sec.~5 and app.~A.1]{arc2025}. Its semantic comparison of candidate
translations should not be simplified into a claim that independent generation
proves equivalence. ARc is a published commercial-system comparator, not an
open-source implementation established by this review.

The requested AAAI paper on trustworthy tax reasoning evaluates combinations of
language models, facts and logic programs, including reviewed statutory rules and
examples. It provides a particularly useful distinction for prototype design:
reasoning with gold rules is a different task from generating the rules. The
human work needed to prepare those rules and examples belongs in the comparison
\citep[Tables~2--3]{jurayj2026,jurayj2025extended}. The reported setting uses
100 numerical tax cases; its application-specific penalty for errors and abstention
cannot be imported as a Hong Kong bank's cost model.

The inspected companion code confirms that extraction and Prolog processing are
separate stages and that execution failure can lead to abstention. LegalMath
should adopt that separation and evaluate facts and rules independently. It
should not quote a tax benchmark's best configuration as evidence that the same
method will reduce private-bank compliance errors. The reading notes also retain
a discrepancy between the paper's narrative attribution of a cost result and
the configuration shown in its table; the proposal relies on the methodological
distinction rather than that ranking \citep{libraryreview}.

The recent \emph{Know Your Limits} preprint studies failures in legal reasoning
and autoformalization, including assumptions that change the intended entailment
task. Its examples strengthen the case for testing entity binding, implicit
assumptions and distinctions between legal interpretation and strict entailment
\citep[secs.~3--5]{limits2026}. It remains under-review evidence. Inconsistencies
in the reported dataset denominators require clarification before quantitative
reuse, and the choice of which ordinary or legal assumptions are admissible is
itself a modeling decision. This proposal borrows its failure taxonomy, not a
general error-rate estimate.

\subsection{Software choices and the role of standards}

The open-source landscape is sufficiently rich that a prototype need not build
every component. It is not sufficiently uniform that selecting one package
settles semantics, review, execution and proof at once. Table~\ref{tab:software}
identifies a practical role for each candidate. These are design recommendations
based on the inspected papers and official documentation; installation,
performance, security and integration have not been tested here.

\begingroup\small
\begin{longtable}{P{0.19\textwidth}P{0.37\textwidth}P{0.34\textwidth}}
\caption{Candidate software and standards, with a deliberately bounded role.}
\label{tab:software}\\
\toprule Candidate & Contribution to LegalMath & Limit or prototype decision \\\midrule
\endfirsthead
\toprule Candidate & Contribution to LegalMath & Limit or prototype decision \\\midrule
\endhead
Catala & Literate definitions, exceptions, executable decisions and verification-condition ideas. & Compare one adapter on the SPI decision fragment; do not assume event-workflow coverage or end-to-end verified compilation. \citep{catalacode}\\
DMN 1.5; Apache KIE / Drools & Standard decision tables and Java-oriented decision execution. & Preserve hit policies, explicit unknown handling and explanations. Add a separate event model. \citep{dmn2024,droolscode}\\
LegalRuleML & Source, authority, interpretation and normative metadata for interchange. & Adopt a small profile; a reasoner and execution semantics are separate components. \citep{legalruleml2021}\\
Stipula; Stipula--KeY & Contract states, permissions, scheduled events and a Java/JML proof experiment. & Research adapter for a small workflow; document event abstraction and loop restrictions. \citep{stipulacode,stipulakeycode}\\
eFLINT & Normative simulation, duties and violations, including actions that occur despite a prohibition. & Borrow concepts and compare a bounded scenario; pin the implementation and semantics version. \citep{eflint2026}\\
Blawx / s(CASP) & Visual rule authoring, explanations and hypothetical reasoning. & Authoring experiment; Blawx describes its current scope as educational and experimental. \citep{blawxcode,slaw2024}\\
Z3; existing Lean route & Counterexamples, satisfiability, selected invariants and checked proof targets. & Specify supported theories and assumptions; preserve timeout and unknown outcomes. \citep{z3code,localrepos}\\
OpenFisca & Versioned computable rules and a tax/benefit microsimulation architecture. & Useful design analogue; the primary problem here includes conduct and workflow duties. \citep{openfiscacode}\\
Open Policy Agent & Rego policies evaluated against structured data; potential deployment adapter. & Its undefined-value behavior must be reconciled with the legal rule semantics. It supplies no natural-language legal interpretation. \citep{opadocs}\\
Accord Project / Cicero & Adjacent work on digital contracts and reusable contract technology. & Monitor for document/template integration; a fixed contract template differs from an amended regulatory corpus. \citep{cicerocode}\\
\bottomrule
\end{longtable}
\endgroup

DMN deserves particular attention because its accessible tables can conceal a
priority choice. Under a Unique hit policy, matching rows must be disjoint; Any
permits overlaps with the same result; First selects by row order. Those policies
are not interchangeable. A default exception tree imported from Catala cannot be
exported as an arbitrary First table without checking that the meanings agree
\citep[sec.~8.2.11]{dmn2024}. Review annotations also need their own source store:
a visual note on a table is not automatically available as runtime provenance.

Licensing and reproducibility are separate from public availability. Catala,
Apache KIE, OpenFisca, Z3 and the other public projects are candidates for reuse,
but the precise pinned version, dependencies and intended distribution must be
reviewed. A paper PDF, a hosted service and a public research repository carry
different rights and maintenance expectations. The Stipula and tax prototypes
should not be described as approved bank components merely because their source
can be inspected. No third-party implementation has been imported into LegalMath
as part of this document work.

\subsection{How to turn the citation network into a development program}

The richest use of this literature is to connect each mechanism to a failure the
bank can recognize. Catala and Logical English address inspectable definitions,
while date semantics and CUTECat address difficult boundaries. LegalRuleML,
s(LAW) and eFLINT address context, unknowns and obligations; Stipula and
process-compliance work address histories. ContractCheck and SMT techniques
address contradictory requirements, and the neurosymbolic papers address the
fallible translation boundary. Each contribution can therefore be tested against
a specific part of the bank's problem.

For each borrowed mechanism, the prototype should retain one source example,
one SFC analogue, the assumptions that changed, and a test that would reveal a
bad transfer. A Stipula event-ordering example becomes a consent-withdrawal trace;
a Catala exception example becomes the fee-discount exception; an ARc ambiguity
case becomes two proposed readings with a concrete client case on which they
disagree. If the mechanism cannot represent the bank case without hidden
conventions, that is evidence to revise the adapter or keep the case under human
assessment. It is not a reason to force the circular into the language.

Citation expansion should continue during development around those unresolved
questions. New papers enter the library with versioned full text and technical
notes, and software enters through a small reproducible example. This preserves
the benefits of a broad research base while keeping the product centered on
reviewable regulatory change.
